\providecommand{\Res}[2]{\qty[round-mode=figures,round-precision=3]{\fpeval{#1}}{#2}}
\providecommand{\Val}[1]{\num[round-mode=figures,round-precision=4]{\fpeval{#1}}}
\subsubsection*{Question a) Accélération}
\begin{enumerate}[label=\arabic*)]
    \item \textbf{Variables connues et recherchées :}
    \begin{itemize}[label= -]
        \item Masse du traîneau : $m = \qty{\v{m}}{kg}$
        \item Tension dans la corde : $T = \qty{\v{T}}{N}$
        \item Angle de la corde : $\theta = \qty{\v{th}}{\degree}$
        \item Forces : $T$ (inclinée), $mg$, $n$
        \item Variable recherchée : $a_x$ (m/s$^2$)
    \end{itemize}
    \item \textbf{Équations pertinentes :}
    \[ \sum F_x = m a_x \implies T\cos\theta = m a_x \]
    \item \textbf{Isolation de la variable recherchée :}
    \[ a_x = \frac{T\cos\theta}{m} \]
    \item \textbf{Résolution numérique :}
    \[ a_x = \frac{\num{\vsi{T}} \cdot \cos(\ang{\v{th}})}{\num{\vsi{m}}} = \Res{\vsi{T}*cosd(\vsi{th})/\vsi{m}}{m/s^2} \]
    \item \textbf{Réponse :} le traîneau accélère à \textbf{\Res{\vsi{T}*cosd(\vsi{th})/\vsi{m}}{m/s^2}}.
\end{enumerate}

\subsubsection*{Question b) Force normale}
\begin{enumerate}[label=\arabic*)]
    \item \textbf{Variables connues et recherchées :}
    \begin{itemize}[label= -]
        \item Masse du traîneau : $m = \qty{\v{m}}{kg}$
        \item Tension dans la corde : $T = \qty{\v{T}}{N}$
        \item Angle de la corde : $\theta = \qty{\v{th}}{\degree}$
        \item $a_y = 0$
        \item Variable recherchée : $n$ (N)
    \end{itemize}
    \item \textbf{Équations pertinentes :}
    \[ \sum F_y = 0 \implies n + T\sin\theta - m g = 0 \]
    \item \textbf{Isolation de la variable recherchée :}
    \[ n = m g - T\sin\theta \]
    \item \textbf{Résolution numérique :}
    \[ n = \num{\vsi{m}} \cdot \num{9.81} - \num{\vsi{T}} \cdot \sin(\ang{\v{th}}) = \Res{\vsi{m}*9.81 - \vsi{T}*sind(\vsi{th})}{N} \]
    \item \textbf{Réponse :} la normale vaut \textbf{\Res{\vsi{m}*9.81 - \vsi{T}*sind(\vsi{th})}{N}}, moins que le poids car la corde soulève en partie le traîneau.
\end{enumerate}

\subsubsection*{Question c) Vitesse}
\begin{enumerate}[label=\arabic*)]
    \item \textbf{Variables connues et recherchées :}
    \begin{itemize}[label= -]
        \item $a_x = \Res{\vsi{T}*cosd(\vsi{th})/\vsi{m}}{m/s^2}$
        \item Durée : $t = \qty{\v{t}}{s}$
        \item $v_0 = 0$
        \item Variable recherchée : $v$ (m/s)
    \end{itemize}
    \item \textbf{Équations pertinentes :}
    \[ v = v_0 + a_x t \]
    \item \textbf{Isolation de la variable recherchée :}
    \[ v = a_x t \]
    \item \textbf{Résolution numérique :}
    \[ v = \Val{\vsi{T}*cosd(\vsi{th})/\vsi{m}} \cdot \num{\vsi{t}} = \Res{\vsi{T}*cosd(\vsi{th})/\vsi{m}*\vsi{t}}{m/s} \]
    \item \textbf{Réponse :} le traîneau atteint \textbf{\Res{\vsi{T}*cosd(\vsi{th})/\vsi{m}*\vsi{t}}{m/s}}.
\end{enumerate}
